% BHU PYQ -- Manually transcribed & error-corrected
\documentclass[11pt,a4paper]{article}

\usepackage[a4paper,top=2.5cm,bottom=2.5cm,left=2.8cm,right=2.8cm,headheight=14pt]{geometry}
\usepackage{amsmath,amssymb}
\usepackage[T1]{fontenc}
\usepackage[utf8]{inputenc}
\usepackage{lmodern}
\usepackage{microtype}
\usepackage{fancyhdr}
\usepackage{enumitem}
\usepackage{hyperref}
\usepackage{lastpage}
\usepackage{parskip}

\hypersetup{colorlinks=true,urlcolor=black,linkcolor=black,
  pdftitle={CHB-04A: Ancillary Chemistry-II -- BHU PYQ},pdfauthor={Banaras Hindu University}}

\pagestyle{fancy}\fancyhf{}
\renewcommand{\headrulewidth}{0.4pt}\renewcommand{\footrulewidth}{0.4pt}
\fancyhead[L]{\small   \textit{Banaras Hindu University}}
\fancyhead[R]{\small   \textit{CHB-04A: Ancillary Chemistry-II}}
\fancyfoot[C]{\small Page  \thepage\ of \pageref{LastPage}}


\newlist{parts}{enumerate}{2}
\setlist[parts,1]{label=(\alph*),leftmargin=*,itemsep=5pt,topsep=3pt}
\setlist[parts,2]{label=(\roman*),leftmargin=*,itemsep=3pt,topsep=2pt}

\newcommand{\pts}[1]{\hfill   \textit{[#1]}}


\begin{document}


\begin{center}
  {\large\bfseries BANARAS HINDU UNIVERSITY}\\[2pt]
  {
\normalsize B.Sc. (Hons.) Semester IV Examination, 2013-14}\\[6pt]
  \rule{0.8\linewidth}{0.5pt}\\[6pt]
  {\large\bfseries Paper No. CHB-04A: Ancillary Chemistry-II (Basic Aspects of Chemistry)}\\[4pt]
  {
\normalsize Time: 3 Hours\hfill Maximum Marks: 70}
\end{center}
\rule{\linewidth}{0.4pt}
\small



\noindent   \textit{Instructions: Answer five questions in total, including Question~1 which is compulsory.}

\normalsize
\bigskip




\noindent   \textbf{Question 1.} Answer all parts of the following: \pts{5 $ \times$ 2 = 10}

\begin{parts}
  \item Discuss in brief the advantages of hydrogen-hydride fuel over conventional fossil fuels.
  \item Write the names of the factors that affect the rate of photosynthesis in plants.
  \item Differentiate between a co-factor and a prosthetic group.
  \item Write the name of any one molecular disease and discuss its reason in brief.
  \item Define the term Respiratory Quotient (RQ).
\end{parts}

\medskip\rule{\linewidth}{0.2pt}




\noindent   \textbf{Question 2.} Differentiate between the following pairs: \pts{7 $ \times$ 2 = 14}

\begin{parts}
  \item Aerobic respiration and anaerobic respiration
  \item Isothermal thermodynamic processes and adiabatic processes
  \item Covalent bonds and coordinate covalent bonds
  \item Extensive thermodynamic properties and intensive properties
\end{parts}

\medskip\rule{\linewidth}{0.2pt}




\noindent   \textbf{Question 3.}

\begin{parts}
  \item Discuss similarity and differences between haemoglobin and myoglobin. Explain the role of haemoglobin in the transport of $ \text{O}_2$ and $ \text{CO}_2$ in human blood. \pts{8}
  \item Discuss the 'light and dark reactions' of photosynthesis in brief. Name the substances which are produced in light reactions and used as inputs in dark reactions. \pts{6}
\end{parts}

\medskip\rule{\linewidth}{0.2pt}




\noindent   \textbf{Question 4.}

\begin{parts}
  \item What do you mean by Nylons? Discuss the reactions involved in the formation of Nylon-6,6 and Nylon-6. Give their uses. \pts{8}
  \item Discuss the energetics (Born-Haber cycle) for the formation of solid $ \text{NaCl}$ from solid sodium and chlorine gas. \pts{6}
\end{parts}

\medskip\rule{\linewidth}{0.2pt}




\noindent   \textbf{Question 5.}

\begin{parts}
  \item Discuss the composition, preparation, and properties of Bakelite and Polyneoprene. \pts{6}
  \item Explain the mechanism of metallic corrosion. How can it be prevented by sacrificial anode protection? \pts{6}
  \item Differentiate between fibrous and globular proteins. \pts{2}
\end{parts}

\vfill
\rule{\linewidth}{0.4pt}

\begin{center}\small
  \href{https://bhu-exam.vercel.app/}{Click here to visit our website}\\[4pt]
  \href{https://me-aryan.vercel.app/}{Click here to visit our contributor}\\[4pt]
  \href{https://quashan-me.vercel.app/}{Click here to visit our contributor}
\end{center}

\end{document}
